\documentclass[12pt,a4paper]{article}

\usepackage[a4paper,margin=2.1cm]{geometry}
\usepackage{fontspec}
\usepackage{polyglossia}
\setmainlanguage{english}
\setmainfont{TeX Gyre Pagella}

\usepackage{xcolor}
\definecolor{maroon}{HTML}{800000}
\definecolor{steel}{HTML}{4682B4}
\definecolor{slate}{HTML}{6A5ACD}
\definecolor{gold}{HTML}{DAA520}
\definecolor{darkgray}{HTML}{333333}

\usepackage{tikz}
\usetikzlibrary{arrows.meta,positioning,shapes.geometric,fit,backgrounds,calc}

\usepackage{booktabs}
\usepackage{tabularx}
\usepackage{array}
\usepackage{enumitem}
\usepackage{tcolorbox}
\usepackage{hyperref}
\usepackage{fancyhdr}

\hypersetup{
  colorlinks=true,
  linkcolor=maroon,
  urlcolor=steel
}

\pagestyle{fancy}
\fancyhf{}
\lhead{\textcolor{maroon}{Pharmacy V13}}
\rhead{\textcolor{steel}{System Architecture and Management}}
\cfoot{\thepage}

\setlength{\parindent}{0pt}
\setlength{\parskip}{6pt}

\newtcolorbox{keybox}{
  colback=maroon!4,
  colframe=maroon!75!black,
  boxrule=0.8pt,
  arc=2mm,
  left=5mm,
  right=5mm,
  top=3mm,
  bottom=3mm
}

\newtcolorbox{notebox}{
  colback=steel!5,
  colframe=steel,
  boxrule=0.7pt,
  arc=2mm,
  left=5mm,
  right=5mm,
  top=3mm,
  bottom=3mm
}

\newtcolorbox{goldbox}{
  colback=gold!8,
  colframe=gold!70!black,
  boxrule=0.7pt,
  arc=2mm,
  left=5mm,
  right=5mm,
  top=3mm,
  bottom=3mm
}

\newcommand{\area}[1]{\textcolor{maroon}{\textbf{#1}}}
\newcommand{\live}{\textcolor{steel}{\textbf{LIVE}}}
\newcommand{\simulated}{\textcolor{slate}{\textbf{SIMULATED DEVELOPMENT DATA}}}

\title{
\textcolor{maroon}{\Huge\textbf{Pharmacy V13}}\\[5pt]
\Large System Architecture, Management and Decision Intelligence
}

\author{
\textbf{Prof. Babiker Malik Osman, PhD}\\
Biomedical Data Scientist \;\textbullet\; System Biologist
}

\date{September 2026}

\begin{document}

\maketitle
\thispagestyle{empty}

\begin{keybox}
\textbf{Pharmacy V13 is a management and decision-intelligence system, not merely an inventory application.}

The system connects operational pharmacy activity with monitoring, quantitative intelligence, evidence-based decision support, management action, reporting, security, testing, and continuous feedback.
\end{keybox}

\section*{Executive Summary}

Pharmacy V13 is an integrated digital pharmacy management platform designed to transform operational activity into structured information, analytical evidence, management decisions, and controlled action.

The central architecture is:

\begin{center}
\begin{tikzpicture}[
node distance=0.55cm,
flow/.style={
rounded corners,
draw=maroon,
very thick,
fill=maroon!5,
minimum width=2.25cm,
minimum height=0.9cm,
align=center,
font=\small\bfseries
},
arrow/.style={-{Latex[length=2.5mm]},thick,gray!70!black}
]
\node[flow] (data) {Operational\\Data};
\node[flow,right=of data] (intel) {Intelligence};
\node[flow,right=of intel] (evidence) {Evidence};
\node[flow,right=of evidence] (decision) {Decision};
\node[flow,right=of decision] (action) {Action};
\node[flow,right=of action] (feedback) {Feedback};

\draw[arrow] (data)--(intel);
\draw[arrow] (intel)--(evidence);
\draw[arrow] (evidence)--(decision);
\draw[arrow] (decision)--(action);
\draw[arrow] (action)--(feedback);
\draw[arrow] (feedback.south) to[out=-90,in=-90] (data.south);
\end{tikzpicture}
\end{center}

The current production environment is deployed through Vercel with PostgreSQL provided by Neon. Authentication, role-based access, protected APIs, operational transactions, reporting, intelligence services, audit mechanisms, and system-health services form one integrated platform.

The verified live production baseline includes:

\begin{itemize}[leftmargin=1cm]
\item 120 medicines;
\item 2 suppliers;
\item 7 recorded sales;
\item 3 purchase records;
\item current Daily Cash revenue of 194;
\item authenticated and role-controlled management access;
\item protected operational APIs;
\item Vercel cloud deployment;
\item Neon PostgreSQL connectivity;
\item REEM V2 decision-intelligence models;
\item reporting and system snapshot functions; and
\item \textbf{45 PASS / 0 FAIL} in the functional verification suite.
\end{itemize}

\begin{goldbox}
\textbf{Important data-integrity principle:}
Live operational data and REEM V2 simulated development data are not interchangeable.

The production database represents actual system transactions. REEM V2 development datasets are used to validate analytical architecture and decision models and must not be presented as live pharmacy observations.
\end{goldbox}

\section*{1. System Purpose}

The purpose of Pharmacy V13 is to provide a coherent management environment in which pharmacy operations, information, analytics, decisions, and technical infrastructure are connected.

The system addresses six fundamental management questions:

\begin{enumerate}[leftmargin=1cm]
\item What is happening?
\item What has happened?
\item What does the data indicate?
\item What deserves attention?
\item What action should management consider?
\item How can the decision and its outcome be documented and reviewed?
\end{enumerate}

This makes Pharmacy V13 a management information and decision-support platform rather than a collection of isolated screens.

\section*{2. Integrated System Architecture}

The system is organized into five functional areas:

\begin{enumerate}[leftmargin=1cm]
\item \area{Control Center} --- see the current state.
\item \area{Operations} --- perform and record daily activity.
\item \area{Decision Intelligence} --- transform data into analytical evidence.
\item \area{Reporting} --- document and communicate results.
\item \area{Engineering} --- provide security, infrastructure, reliability, testing, and maintenance.
\end{enumerate}

\begin{center}
\begin{tikzpicture}[
box/.style={
rounded corners,
draw=#1,
very thick,
fill=#1!7,
minimum width=3.1cm,
minimum height=1.15cm,
align=center,
font=\small\bfseries
},
arrow/.style={-{Latex[length=3mm]},thick,gray!70!black}
]

\node[box=maroon] (control) {CONTROL CENTER\\See};
\node[box=steel,right=0.75cm of control] (operations) {OPERATIONS\\Do};
\node[box=slate,right=0.75cm of operations] (decision) {DECISION\\Understand};
\node[box=gold,right=0.75cm of decision] (report) {REPORTING\\Communicate};

\node[box=maroon,below=1.25cm of operations,minimum width=4.2cm]
(engineering) {ENGINEERING\\Build, Protect \& Maintain};

\draw[arrow] (operations)--(control);
\draw[arrow] (operations)--(decision);
\draw[arrow] (decision)--(report);
\draw[arrow] (engineering)--(operations);
\draw[arrow] (engineering)--(decision);
\draw[arrow] (engineering)--(report);

\end{tikzpicture}
\end{center}

These are different views of one integrated system, not independent applications.

\section*{3. Cloud and Technical Architecture}

The current production architecture separates the management interface, API services, database, authentication, and analytical intelligence.

\begin{center}
\begin{tikzpicture}[
layer/.style={
rounded corners,
draw=maroon,
thick,
minimum width=11cm,
minimum height=0.82cm,
align=center,
fill=maroon!5
},
arr/.style={-{Latex[length=2.5mm]},thick,gray}
]

\node[layer] (web) {Web Application / Management Interface};
\node[layer,below=0.35cm of web] (api) {Vercel API Services / Protected Endpoints};
\node[layer,below=0.35cm of api] (auth) {Authentication / JWT / Role-Based Access};
\node[layer,below=0.35cm of auth] (db) {Neon PostgreSQL / Operational Data};
\node[layer,below=0.35cm of db] (reem) {REEM V2 / Analytical Evidence / AI Interpretation};

\draw[arr] (web)--(api);
\draw[arr] (api)--(auth);
\draw[arr] (auth)--(db);
\draw[arr] (db)--(reem);

\end{tikzpicture}
\end{center}

\subsection*{Production Components}

\begin{itemize}[leftmargin=1cm]
\item \textbf{Frontend:} browser-based Pharmacy V13 management interface.
\item \textbf{Cloud platform:} Vercel production deployment.
\item \textbf{Database:} Neon PostgreSQL.
\item \textbf{API layer:} protected Vercel API services.
\item \textbf{Authentication:} application-managed authentication using secure password hashing and JWT sessions.
\item \textbf{Authorization:} role-based access control.
\item \textbf{Analytics:} REEM V2 decision-intelligence layer.
\item \textbf{Reporting:} controlled system reports, snapshots, and analytical outputs.
\end{itemize}

\section*{4. Authentication and Role-Based Access}

Security is part of the system architecture rather than an optional interface feature.

Pharmacy V13 uses application-specific authentication integrated with the production database.

The principal roles are:

\begin{center}
\begin{tabularx}{0.9\textwidth}{lX}
\toprule
\textbf{Role} & \textbf{Management responsibility}\\
\midrule
Administrator & Full system administration and privileged operations.\\
Manager & Operational management and controlled transaction functions.\\
Analyst & Analysis, intelligence, reporting, and evidence review.\\
Viewer & Read-only access to authorized information.\\
\bottomrule
\end{tabularx}
\end{center}

Protected APIs require authentication. The health endpoint remains publicly available for service monitoring.

Write operations such as sales and inventory stock transactions are restricted to authorized management roles.

\begin{notebox}
\textbf{Security principle:}
Authentication establishes who the user is; authorization determines what that user is permitted to do.
\end{notebox}

\section*{5. Control Center --- See}

The Control Center provides the management view of the current operational state.

Typical indicators include:

\begin{itemize}[leftmargin=1cm]
\item medicine count;
\item sales and revenue;
\item inventory quantity and value;
\item stock alerts;
\item expiry status;
\item supplier and procurement status;
\item decision-intelligence status;
\item daily cash;
\item system health; and
\item management snapshots.
\end{itemize}

The Control Center answers:

\begin{center}
\textbf{What is happening in the pharmacy now?}
\end{center}

\section*{6. Pharmacy Operations --- Do}

The operational layer records the events that create the pharmacy's digital state.

The main operational entities are:

\begin{itemize}[leftmargin=1cm]
\item medicines;
\item suppliers;
\item purchases;
\item inventory;
\item sales;
\item stock-in transactions;
\item stock-out transactions;
\item cash activity; and
\item operational alerts.
\end{itemize}

\begin{center}
\begin{tikzpicture}[
node distance=0.7cm,
flow/.style={
rounded corners,
draw=steel,
thick,
fill=steel!7,
minimum width=2.35cm,
minimum height=0.85cm,
align=center,
font=\small\bfseries
},
a/.style={-{Latex[length=2.5mm]},thick,steel}
]

\node[flow] (supplier) {Supplier};
\node[flow,right=of supplier] (purchase) {Purchase};
\node[flow,right=of purchase] (inventory) {Inventory};
\node[flow,right=of inventory] (sale) {Sale};
\node[flow,right=of sale] (cash) {Cash};

\draw[a] (supplier)--(purchase);
\draw[a] (purchase)--(inventory);
\draw[a] (inventory)--(sale);
\draw[a] (sale)--(cash);

\end{tikzpicture}
\end{center}

Every transaction contributes to the evidence available for management and analytical decision-making.

\section*{7. Daily Cash}

Daily Cash provides a concise financial view of recorded sales.

The current production checkpoint reports:

\begin{center}
\begin{tabular}{lr}
\toprule
\textbf{Indicator} & \textbf{Current value}\\
\midrule
Recorded sales & 7\\
Revenue & 194\\
\bottomrule
\end{tabular}
\end{center}

The purpose of Daily Cash is not to replace formal accounting. It provides an operational management view derived from system transactions.

\section*{8. REEM V2 Decision Intelligence}

REEM V2 is the analytical decision-intelligence layer of Pharmacy V13.

Its role is to transform structured operational or development data into analytical evidence that can be interpreted by management and AI.

The principal models are:

\begin{enumerate}[leftmargin=1cm]
\item Demand Forecasting
\item ABC Analysis
\item Expiry Risk
\item Stock Coverage
\item Supplier Analysis
\item Reorder Point
\item Safety Stock
\item Economic Order Quantity
\end{enumerate}

Additional risk analysis includes Monte Carlo simulation where appropriate.

\begin{keybox}
\textbf{REEM V2 calculates evidence; AI interprets the evidence.}

REEM V2 is responsible for quantitative analytical evidence. AI provides a management interpretation of that evidence. The two functions should remain conceptually distinguishable.
\end{keybox}

\section*{9. Forecasting}

Demand forecasting estimates future demand from available historical demand information.

The purpose is not simply to produce a number. Forecasting provides a quantitative basis for:

\begin{itemize}[leftmargin=1cm]
\item procurement planning;
\item stock planning;
\item replenishment decisions;
\item capacity planning; and
\item identification of changing demand patterns.
\end{itemize}

Forecasts should always be interpreted together with data quality, observation history, and uncertainty.

\section*{10. ABC Analysis}

ABC analysis prioritizes medicines according to their relative contribution to the selected value measure.

Its management purpose is to answer:

\begin{center}
\textbf{Which medicines deserve the greatest management attention?}
\end{center}

ABC classification can support differentiated control policies, monitoring frequency, procurement attention, and resource allocation.

It is therefore a prioritization tool rather than a simple classification table.

\section*{11. Reorder Point and Safety Stock}

Reorder Point determines when replenishment should be initiated.

Safety Stock provides protection against uncertainty in demand and replenishment lead time.

Conceptually:

\begin{center}
\textbf{Reorder Point = Expected Lead-Time Demand + Safety Stock}
\end{center}

The models therefore connect demand evidence with operational replenishment decisions.

\begin{notebox}
A replenishment signal should be interpreted in relation to current stock, expected demand, lead time, uncertainty, and management policy.
\end{notebox}

\section*{12. Expiry Risk}

Expiry analysis identifies medicines and batches that require management attention because of remaining shelf life.

Expiry intelligence can support:

\begin{itemize}[leftmargin=1cm]
\item prioritization of stock;
\item FEFO-oriented management;
\item reduction of avoidable losses;
\item procurement review; and
\item early intervention before expiry.
\end{itemize}

The analytical distinction between batch-level observations and medicine-level summaries must be preserved when interpreting expiry results.

\section*{13. Economic Order Quantity}

EOQ provides a classical analytical framework for balancing ordering and holding costs under its underlying assumptions.

The conceptual model is:

\begin{center}
\textbf{Optimal Order Quantity balances ordering cost and holding cost.}
\end{center}

EOQ is therefore a decision-support model, not an automatic purchasing command.

Management must consider supplier constraints, minimum order quantities, expiry, cash availability, demand uncertainty, and practical procurement conditions.

\section*{14. Supplier Intelligence}

Supplier intelligence connects procurement records with supplier-level performance information.

Relevant evidence may include:

\begin{itemize}[leftmargin=1cm]
\item purchasing activity;
\item supply reliability;
\item delivery performance;
\item procurement history;
\item medicine availability; and
\item risk indicators.
\end{itemize}

The objective is to support supplier evaluation and procurement strategy rather than merely maintain a supplier directory.

\section*{15. Monte Carlo Risk Analysis}

Monte Carlo analysis provides a mechanism for exploring uncertainty by repeatedly simulating possible outcomes.

In Pharmacy V13, this approach can support questions such as:

\begin{itemize}[leftmargin=1cm]
\item How sensitive is stock availability to demand variation?
\item How much safety stock may be required under uncertainty?
\item What is the probability of a stock-out?
\item How robust is a replenishment decision?
\end{itemize}

Simulation outputs are evidence about possible outcomes, not predictions of certainty.

\section*{16. AI Management Interpretation}

AI is positioned above the analytical evidence rather than replacing the underlying models.

The intended chain is:

\begin{center}
\begin{tikzpicture}[
node distance=0.55cm,
flow/.style={
rounded corners,
draw=slate,
very thick,
fill=slate!6,
minimum width=3cm,
minimum height=0.9cm,
align=center,
font=\small\bfseries
},
a/.style={-{Latex[length=2.5mm]},thick,slate}
]

\node[flow] (data) {Operational /\\Analytical Data};
\node[flow,right=of data] (model) {REEM V2\\Models};
\node[flow,right=of model] (evidence) {Evidence};
\node[flow,right=of evidence] (ai) {AI\\Interpretation};
\node[flow,right=of ai] (management) {Management\\Decision};

\draw[a] (data)--(model);
\draw[a] (model)--(evidence);
\draw[a] (evidence)--(ai);
\draw[a] (ai)--(management);

\end{tikzpicture}
\end{center}

This architecture improves transparency because management recommendations can be traced back to analytical evidence.

\section*{17. Reporting and Reproducibility}

The Report Center provides controlled documentation of the system and its analytical outputs.

Reporting functions include:

\begin{itemize}[leftmargin=1cm]
\item Pharmacy V13 system documentation;
\item REEM V2 analytical documentation;
\item conceptual architecture documentation;
\item management reports;
\item SAS-oriented analytical outputs;
\item system snapshots;
\item system-health information; and
\item reproducible source documents.
\end{itemize}

A report is therefore not simply a printout.

\begin{keybox}
\textbf{A report is a controlled representation of what the system knows at a particular point in time.}
\end{keybox}

The LaTeX source provides a reproducible documentation layer for the principal system report.

\section*{18. Security, Audit and Data Integrity}

Pharmacy V13 treats security and data integrity as architectural requirements.

The security architecture includes:

\begin{itemize}[leftmargin=1cm]
\item authenticated access;
\item role-based authorization;
\item protected operational APIs;
\item secure password hashing;
\item signed session tokens;
\item login-attempt recording;
\item application audit records;
\item controlled write operations; and
\item public health monitoring separated from protected business APIs.
\end{itemize}

Data integrity also requires a clear distinction between:

\begin{enumerate}[leftmargin=1cm]
\item live operational data;
\item analytical transformations;
\item simulated development datasets;
\item model outputs; and
\item AI interpretations.
\end{enumerate}

\begin{goldbox}
\textbf{No simulated REEM V2 development result should be represented as live pharmacy evidence.}
\end{goldbox}

\section*{19. Testing and Verification}

The Pharmacy V13 functional verification suite currently reports:

\begin{center}
\begin{tikzpicture}
\node[
rounded corners,
draw=steel,
very thick,
fill=steel!7,
minimum width=7cm,
minimum height=1.4cm,
align=center,
font=\Large\bfseries
] {45 PASS \qquad 0 FAIL};
\end{tikzpicture}
\end{center}

The test suite provides a controlled baseline for subsequent development.

Testing is treated as part of the engineering architecture rather than an activity performed only before release.

\begin{notebox}
\textbf{Engineering principle:}
A change is not considered complete merely because the application appears to work. The existing verification baseline must remain intact.
\end{notebox}

\section*{20. Current Production Status}

The current Pharmacy V13 production environment is deployed at:

\begin{center}
\url{https://pharmacy-v13.vercel.app/}
\end{center}

The production health service confirms:

\begin{itemize}[leftmargin=1cm]
\item service status: healthy;
\item production API: operational;
\item Neon PostgreSQL connection: active;
\item database: \texttt{neondb};
\item medicines: 120;
\item sales: 7;
\item suppliers: 2;
\item purchases: 3.
\end{itemize}

The production application therefore represents a functioning cloud management platform rather than a purely local prototype.

\section*{21. Management Decision Cycle}

The complete management cycle can be represented as:

\begin{center}
\begin{tikzpicture}[
node distance=0.55cm,
step/.style={
rounded corners,
draw=maroon,
thick,
fill=maroon!5,
minimum width=2.5cm,
minimum height=0.85cm,
align=center,
font=\small\bfseries
},
a/.style={-{Latex[length=2.5mm]},thick,maroon}
]

\node[step] (observe) {Observe};
\node[step,right=of observe] (record) {Record};
\node[step,right=of record] (analyze) {Analyze};
\node[step,right=of analyze] (decide) {Decide};
\node[step,right=of decide] (act) {Act};
\node[step,right=of act] (review) {Review};

\draw[a] (observe)--(record);
\draw[a] (record)--(analyze);
\draw[a] (analyze)--(decide);
\draw[a] (decide)--(act);
\draw[a] (act)--(review);
\draw[a] (review.south) to[out=-90,in=-90] (observe.south);

\end{tikzpicture}
\end{center}

This cycle is the conceptual core of Pharmacy V13.

The system observes operational reality, records it, analyzes evidence, supports management decisions, records action, and uses subsequent information for review and improvement.

\section*{22. One Integrated Pharmacy Scenario}

Consider a medicine approaching its replenishment threshold.

\begin{enumerate}[leftmargin=1cm]
\item A purchase is received from a supplier.
\item The purchase is recorded.
\item Inventory is updated.
\item The Control Center reflects the new operational state.
\item Sales subsequently reduce available stock.
\item Demand evidence is updated.
\item REEM V2 evaluates demand, coverage, reorder point, safety stock, expiry, and related risk.
\item The analytical layer produces evidence.
\item AI interprets the evidence for management.
\item Management decides whether and how to replenish.
\item The resulting action becomes a new operational event.
\item Reporting and snapshots preserve the relevant state for review.
\end{enumerate}

Thus the system closes the loop between operations and decision-making.

\section*{23. Live Data versus Development Intelligence}

This distinction is essential to the scientific integrity of Pharmacy V13.

\begin{center}
\begin{tabularx}{0.92\textwidth}{lXX}
\toprule
\textbf{Layer} & \textbf{Purpose} & \textbf{Status}\\
\midrule
Operational PostgreSQL & Real system transactions & \live\\
Authentication and audit & Access and accountability & \live\\
Production APIs & Operational services & \live\\
Control Center & Management interface & \live\\
REEM V2 datasets & Model development and validation & \simulated\\
REEM V2 models & Decision-intelligence architecture & Development / analytical\\
AI interpretation & Evidence interpretation & Development / analytical\\
\bottomrule
\end{tabularx}
\end{center}

The distinction allows the analytical architecture to evolve without falsely claiming that simulated development evidence represents live pharmacy activity.

\section*{24. Engineering Philosophy}

Pharmacy V13 follows a deliberately modular architecture.

The system should evolve through controlled increments rather than uncontrolled rewrites.

The principal engineering principles are:

\begin{itemize}[leftmargin=1cm]
\item preserve working functionality;
\item make the smallest necessary change;
\item protect operational data;
\item preserve authentication boundaries;
\item maintain test coverage;
\item separate operational truth from analytical simulation;
\item keep evidence traceable;
\item document important changes;
\item maintain reproducibility; and
\item deploy only verified functionality.
\end{itemize}

\section*{25. Strategic View}

Pharmacy V13 can be understood as a small-scale example of a broader digital health and biomedical information architecture.

Its important contribution is not the individual inventory screen or the individual analytical model.

Its value lies in connecting:

\begin{center}
\textbf{Data}
$\rightarrow$
\textbf{Information}
$\rightarrow$
\textbf{Evidence}
$\rightarrow$
\textbf{Decision}
$\rightarrow$
\textbf{Action}
$\rightarrow$
\textbf{Learning}
\end{center}

This architecture can support future extensions in supply-chain intelligence, pharmacological information management, biomedical analytics, forecasting, quality assurance, and AI-assisted decision support.

\section*{Conclusion}

Pharmacy V13 is a unified management, analytical, and engineering platform.

Its architecture connects daily operational transactions with cloud infrastructure, authentication, PostgreSQL data, decision intelligence, AI interpretation, management action, reporting, and feedback.

The fundamental principle is:

\begin{keybox}
\textbf{Operations create data.}\\
\textbf{The system organizes the data.}\\
\textbf{REEM V2 creates analytical evidence.}\\
\textbf{AI interprets the evidence.}\\
\textbf{Management makes decisions.}\\
\textbf{Actions create new operational data.}\\
\textbf{Reporting preserves and communicates the result.}
\end{keybox}

This closes the management loop and establishes Pharmacy V13 as a practical digital decision-support architecture rather than merely an inventory application.

\vfill

\begin{center}
\textcolor{maroon}{\rule{0.75\textwidth}{0.6pt}}\\[4pt]
\textbf{Pharmacy V13}\\
System Architecture, Management and Decision Intelligence\\
September 2026
\end{center}

\end{document}

